% !TEX program = xelatex
% ============================================================
%  عرض تقديمي تجاري — Sales Pitch Template
%  قالب عرض تقديمي لعرض منتج/خدمة على المستثمرين أو العملاء
%  Compile with: xelatex main.tex (run twice)
% ============================================================
%  Features:
%    - 16:9 widescreen, clean default beamer theme with accent color
%    - Arabic RTL via polyglossia + Amiri font
%    - Problem, solution, demo, market, business model, traction,
%      competition, team, ask, contact slides
%    - Market and competition tables with placeholder numbers
%    - English terms wrapped with \textEN{...}
% ============================================================

\documentclass[aspectratio=169,11pt]{beamer}

% ---- Arabic / RTL language support ----
\usepackage{polyglossia}
\setdefaultlanguage[calendar=gregorian,hijricorrection=1]{arabic}
\setotherlanguage[variant=british]{english}

% ---- Fonts (beamer uses sans by default, so set sans to Amiri too) ----
\usepackage[no-math]{fontspec}
\setmainfont[Scale=1]{NotoKufiArabic-Regular.ttf}
\setsansfont[Scale=1]{NotoKufiArabic-Regular.ttf}
\newfontfamily\arabicfont[Script=Arabic,Scale=0.95]{NotoKufiArabic-Regular.ttf}
\newfontfamily\arabicfontsf[Script=Arabic,Scale=0.95]{NotoKufiArabic-Regular.ttf}
\newfontfamily\englishfont[Scale=0.95]{TeX Gyre Termes}
\let\textEN\textenglish

% ---- Core packages ----
\usepackage{graphicx}
\usepackage{booktabs}
\usepackage{tabularx}
\usepackage{tikz}
\usepackage{amsmath}
\usepackage{amssymb}

% ---- Theme & appearance ----
\usetheme{default}
\setbeamertemplate{navigation symbols}{}
\setbeamersize{text margin left=10mm, text margin right=10mm}

% ---- Accent color ----
\definecolor{accent}{HTML}{B71C1C}
\setbeamercolor{frametitle}{fg=accent}
\setbeamercolor{title}{fg=accent}
\setbeamercolor{alerted text}{fg=accent}
\setbeamercolor{structure}{fg=accent}
\setbeamertemplate{frametitle}[default][center]

% ---- Metadata ----
% TODO: استبدل باسم الشركة ومقدم العرض الفعلي
\title{اسم الشركة أو المنتج}
\subtitle{شعار مختصر يصف القيمة}
\author{اسم المقدّم}
\institute{الشركة}
\date{شهر سنة}

\begin{document}

% ============================================================
% TITLE SLIDE
% ============================================================
\begin{frame}[plain]
\titlepage
\end{frame}

% ============================================================
% THE PROBLEM
% ============================================================
\begin{frame}{المشكلة}
% TODO: صف المشكلة التي يحلها المنتج
\begin{itemize}
  \item المشكلة التي يواجهها العملاء المستهدفون.
  \item حجم المعاناة وتكلفتها \textEN{(cost of the problem)}.
  \item لماذا لم تُحل المشكلة بشكل كافٍ حتى الآن.
\end{itemize}
\end{frame}

% ============================================================
% THE SOLUTION
% ============================================================
\begin{frame}{الحل}
% TODO: قدّم الحل بإيجاز
\begin{itemize}
  \item وصف مختصر للمنتج/الخدمة.
  \item القيمة الأساسية للعميل \textEN{(value proposition)}.
  \item ما الذي يجعل الحل مميزاً.
\end{itemize}
\end{frame}

% ============================================================
% PRODUCT DEMO
% ============================================================
\begin{frame}{عرض المنتج}
% TODO: اذكر الميزات الرئيسية
\begin{itemize}
  \item الميزة الأولى وفائدتها للعميل.
  \item الميزة الثانية \textEN{(feature)}.
  \item الميزة الثالثة مع لقطة شاشة إن وجدت.
\end{itemize}
\end{frame}

% ============================================================
% MARKET OPPORTUNITY
% ============================================================
\begin{frame}{فرصة السوق}
% TODO: استبدل الأرقام بالقيم الفعلية
\begin{table}[h]
\centering
\begin{tabularx}{0.85\textwidth}{lX}
\toprule
المؤشر & القيمة \\
\midrule
السوق الكلي \textEN{(TAM)} & 5 مليار دولار \\
السوق المتاح \textEN{(SAM)} & 1.2 مليار دولار \\
السوق المستهدف \textEN{(SOM)} & 150 مليون دولار \\
معدل النمو السنوي & 12\% \\
\bottomrule
\end{tabularx}
\end{table}
\end{frame}

% ============================================================
% BUSINESS MODEL
% ============================================================
\begin{frame}{نموذج العمل}
% TODO: وضّح كيف تتحقق الإيرادات
\begin{itemize}
  \item مصادر الإيراد الرئيسية \textEN{(revenue streams)}.
  \item نموذج التسعير.
  \item قنوات البيع والتوزيع.
\end{itemize}
\end{frame}

% ============================================================
% TRACTION / METRICS
% ============================================================
\begin{frame}{الإنجازات والمؤشرات}
% TODO: اذكر الأرقام الفعلية
\begin{itemize}
  \item عدد العملاء الحاليين.
  \item الإيراد الشهري المتكرر \textEN{(MRR)}.
  \item معدل النمو شهرياً.
  \item شراكات أو عقود موقعة.
\end{itemize}
\end{frame}

% ============================================================
% COMPETITION
% ============================================================
\begin{frame}{المنافسة}
% TODO: عدّل الجدول حسب المنافسين الفعليين
\begin{table}[h]
\centering
\begin{tabularx}{0.9\textwidth}{lccc}
\toprule
المعيار & منافس أ & منافس ب & نحن \\
\midrule
السعر & مرتفع & متوسط & منخفض \\
سهولة الاستخدام & متوسطة & منخفضة & عالية \\
التغطية & محدودة & واسعة & واسعة \\
\bottomrule
\end{tabularx}
\end{table}
\end{frame}

% ============================================================
% TEAM
% ============================================================
\begin{frame}{الفريق}
% TODO: أضف أسماء ومسؤوليات الفريق الفعلية
\begin{itemize}
  \item اسم المؤسس — المنصب والخلفية.
  \item اسم العضو الثاني — المنصب والخبرة.
  \item اسم العضو الثالث \textEN{(role)}.
\end{itemize}
\end{frame}

% ============================================================
% THE ASK
% ============================================================
\begin{frame}{الطلب}
% TODO: حدد المبلغ والخطوات التالية
\begin{itemize}
  \item المبلغ المطلوب من الاستثمار.
  \item استخدامات التمويل \textEN{(use of funds)}.
  \item الخطوات والمعالم القادمة خلال 12 شهراً.
\end{itemize}
\end{frame}

% ============================================================
% CONTACT
% ============================================================
\begin{frame}[plain]
\begin{center}
\vspace{1.2cm}
{\Huge\bfseries\color{accent} شكراً لكم}\\[1cm]
{\Large تواصلوا معنا}\\[0.8cm]
{\large اسم المقدّم}\\
{\normalsize البريد: \textEN{contact@company.com}}\\
{\normalsize الموقع: \textEN{www.company.com}}
\end{center}
\end{frame}

\end{document}
